% ==================================================================
% id:       app-005
% topic:    שימושי הנגזרת
% subtopic: הוכחת אי־שוויונות
% solution: full
% ==================================================================
הוכיחו כי לכל \(x \geq 1\) מתקיים $\ds{x \ln x \leq \frac{1}{2}(x^2 - 1)}$.
\hint{העבירו אגפים והגדירו פונקציית עזר. בדקו את ערכה ב-$x=1$ ואת סימן הנגזרת שלה (ייתכן שתצטרכו לגזור פעמיים).}
\begin{solbox}
תהא
$f(x)=\frac{1}{2}(x^2-1)-x\ln x$.
מתקיים
$f(1)=0$
ו-
$$.f'(x)=x-1-\ln x$$
מתקיים
$$f''(x)=1-\frac{1}{x}>0$$
לכל
$x>1$,
לכן
$f'$
מונוטונית עולה ב-
$[1,\infty)$,
בפרט
$f'(x)>f'(1)=0$
לכל
$x>1$.
מכאן נסיק כי בתחום
$[1,\infty)$
הפונקציה
$f$
עולה,
בפרט
$f(x)>f(1)=0$,
מכאן הטענה.
\iffalse
נציג פתרון חלופי.
נזכר בתוצאה של סעיף ג' של שאלה \qref{app-003}, לפיה
$f'(x)>0$
(אכן, תקבלו את התוצאה עבור
$b=x$
ו-
$a=1$).
מה קיבלנו? הערך של פונקציה
$f$
ב-
$x=1$
הוא
$0$
ונגזרתה חיובית, כלומר,
$f$
עולה בתחום
$(1,\infty)$,
בפרט
$$f(x)>f(1)=0$$
לכל
$x>1$,
זאת אומרת
$$x\ln x<\frac{1}{2}(x^2-1)$$
לכל
$x>1$,
מכיוון שעבור
$x=0$
מתקבל שיוויון, נסיק כי
$$x\ln x\leq\frac{1}{2}(x^2-1)$$
לכל
$x\geq1$,
כנדרש.
\fi
\end{solbox}
